\section{Continued Fractions and the Gap Structure}

The continued fraction expansion of $\alpha$ is not just a tool in the proof---it
is the \emph{exact} dictionary between arithmetic and geometry in the three-distance theorem.

\subsection*{Continued Fraction Recap}

Every irrational $\alpha \in (0,1)$ has a unique continued fraction expansion
\[
  \alpha = \cfrac{1}{a_1 + \cfrac{1}{a_2 + \cfrac{1}{a_3 + \cdots}}} = [0; a_1, a_2, a_3, \ldots]
\]
with partial quotients $a_k \in \mathbb{N}$.  The \emph{convergents} $p_k/q_k = [0; a_1, \ldots, a_k]$
are the best rational approximations to $\alpha$ and satisfy
\[
  |q_k \alpha - p_k| = \|q_k \alpha\| \;<\; \frac{1}{q_{k+1}}.
\]

\subsection*{Denominators as Transition Points}

\begin{theorem}
  The gap structure of $S_n(\alpha)$ transitions from two lengths to three lengths
  (and back to two) precisely at the \emph{convergent denominators} $q_k$.
  Concretely:
  \begin{itemize}
    \item For $n = q_k$: exactly \textbf{two} distinct gap lengths.
    \item For $q_k < n < q_{k+1}$: exactly \textbf{three} distinct gap lengths,
          with the third being $L_k - \ell_k = \|q_{k-1}\alpha\| - (n - q_k)\|q_k\alpha\|$... adjusted.
    \item At $n = q_{k+1}$: back to \textbf{two} gap lengths.
  \end{itemize}
\end{theorem}

\subsection*{The Role of Partial Quotients}

Large partial quotients $a_k$ mean $q_{k+1} = a_k q_k + q_{k-1}$ is much larger than $q_k$.
During the long interval $n \in (q_k, q_{k+1})$, the three-gap structure persists,
with the longest gap $L_k$ being split $a_k$ times before the new long gap $L_{k+1}$
takes over.  Thus:
\begin{itemize}
  \item \textbf{Large $a_k$}: the three-gap regime lasts a long time; $\alpha$ is
        ``well approximated'' by $p_k/q_k$.
  \item \textbf{All $a_k = 1$}: this is the golden ratio $\varphi = [0;1,1,1,\ldots]$,
        where transitions happen as often as possible (Fibonacci denominators),
        and at each transition the two gap lengths are in ratio $1:\varphi$.
\end{itemize}

\subsection*{Sturmian Sequences}

The three-distance theorem has a combinatorial shadow: encode each gap as a letter
($\mathtt{S}$ for short, $\mathtt{L}$ for long).  Reading the gaps in order around
the circle gives a \emph{Sturmian word}---an infinite binary sequence of minimal
complexity that is aperiodic yet highly structured.  The three-distance property
is equivalent to the statement that Sturmian words have \emph{exactly two distinct
letter frequencies}, differing by $1/n$ at each stage.

\subsection*{The Farey Graph Connection}

The denominators $q_k$ are the vertices of the Stern--Brocot tree (equivalently,
the Farey graph) visited by the continued fraction algorithm.  The gap lengths
$\|q_k \alpha\|$ and $\|q_{k-1}\alpha\|$ are the two \emph{Farey neighbours} of
$\alpha$ at resolution $q_k$.  The three-distance theorem can thus be read as a
statement about the Farey graph: at every resolution, $\alpha$ sits in a Farey
interval whose endpoints provide the two (or three) gap lengths.
